\chapter{Introduction}
\label{ch:introduction}

\section{Background and Motivation}
Modern operating systems generate a continuous stream of low-level events such as
process creation, file access, and network connections. These events capture how a
system behaves over time and can be used to identify malicious activity. Traditional
signature-based intrusion detection systems are effective for known threats but are
unable to recognize new attack patterns. As attack techniques evolve rapidly, there
is a strong need for zero-day detection methods that do not depend on labeled attack
data.

Graph-based representations provide an effective way to capture relationships among
system entities (processes, files, and sockets). In parallel, self-supervised learning
has emerged as a promising approach for modeling normal behavior without explicit
labels. This project integrates these directions to build a practical zero-day attack
detection pipeline that learns from normal activity and flags deviations as anomalies.

\section{Challenges}
Despite the promise of data-driven security, several challenges arise in real-time
host-based detection.

\subsection{Lack of Labeled Attack Data}
Collecting labeled attack traces for every new threat is expensive and often
impossible. A practical system must operate without attack labels while still
achieving useful detection performance.

\subsection{Heterogeneous System Behavior}
System activity involves multiple entity types and interactions. Modeling only
individual events ignores the structural relationships that distinguish normal
workflows from attacks.

\subsection{Dynamic and Noisy Environments}
Normal system behavior varies over time depending on user activity, installed
software, and background services. This variability introduces noise and increases
the risk of false positives.

\section{Problem Statement}
The central problem is to design a host-based detection system that can identify
zero-day attacks without relying on labeled attack data. The system must capture
complex behavioral relationships among system entities, learn normal patterns from
real-time monitoring, and detect anomalous behavior while remaining lightweight
enough for practical deployment.

\section{Proposed Solution}
This project proposes a self-supervised behavioral graph-based detection system.
It monitors real system activity, constructs heterogeneous graphs, and trains a
graph autoencoder to learn normal behavior. During testing, a hybrid anomaly
detector combines reconstruction error with structural features to detect suspicious
activity. Simulated attack scenarios are used to evaluate detection performance in
a controlled environment.

\subsection{Graph-Based Behavioral Modeling}
Events are converted into graphs with PROCESS, FILE, and SOCKET nodes and edges
representing interactions such as executes, reads, writes, spawns, and connects.
This representation captures the context of system activity rather than isolated
event counts.

\subsection{Self-Supervised Graph Autoencoding}
A graph autoencoder learns latent representations of normal graphs and attempts to
reconstruct their structure. Graphs that deviate significantly from learned patterns
produce high reconstruction errors and are flagged as anomalies.

\subsection{Hybrid Anomaly Scoring}
Anomaly scores are computed using both reconstruction error and structural
statistics such as node-type ratios, degree distributions, and edge density. This
hybrid approach improves robustness compared to using reconstruction error alone.

\section{Objectives}
The objectives of this project are:
\begin{enumerate}
	\item Collect real-time system events for processes, files, and network
	connections.
	\item Construct heterogeneous temporal graphs capturing behavioral
	relationships among system entities.
	\item Train a self-supervised graph autoencoder to learn normal behavior
	without attack labels.
	\item Detect anomalies using combined reconstruction and structural features.
	\item Simulate representative attacks to evaluate detection performance.
	\item Generate visualizations and logs to support analysis and reporting.
\end{enumerate}

\section{Scope of the Project}
\textbf{In Scope:} real-time host monitoring, behavioral graph construction,
self-supervised graph autoencoding, hybrid anomaly detection, simulated attack
scenarios (reverse shell, privilege escalation, data exfiltration), performance
metrics, and visualization of results.

\textbf{Out of Scope:} kernel-level instrumentation, enterprise-scale distributed
deployment, full MITRE ATT\&CK mapping, and automated incident response. These
can be addressed in later phases as extensions.

\section{Organization of the Report}
The chapters in this report are organized as follows:

\vspace{0.6em}
CHAPTER 1 INTRODUCTION presents the problem, challenges, and objectives of the
project.

\vspace{0.6em}
CHAPTER 2 LITERATURE SURVEY reviews related work in self-supervised anomaly
detection, graph-based intrusion detection, and continual learning.

\vspace{0.6em}
CHAPTER 3 SYSTEM DESIGN describes the system architecture and module design of
the proposed solution.

\vspace{0.6em}
Subsequent chapters cover implementation, experiments, results, and conclusions.
